\documentclass[10pt,a4paper]{amsart}
\usepackage[margin=19mm]{geometry}
\usepackage{amsmath,amssymb,mathtools}
\usepackage[scr=rsfs,cal=euler]{mathalfa}
\usepackage{xfrac}
\usepackage[unicode,hidelinks,bookmarks=false]{hyperref}
\newcommand{\ie}{i.e.\ }
\newcommand{\cf}{cf.\ }
\newcommand{\resp}{respectively\ }
\newcommand{\Subsection}{\subsection}
\providecommand{\nobreakdash}{\nobreak\hbox{-}}
\setlength{\emergencystretch}{2em}
\multlinegap0pt
\usepackage{bm}
\usepackage{bbm}
\usepackage{dsfont}
\usepackage{xspace}
\usepackage{cancel}
\usepackage{mathtools}
\usepackage{amsthm}
\makeatletter
\@ifundefined{lemma}{\newtheorem{lemma}{Lemma}}{}
\@ifundefined{theorem}{\newtheorem{theorem}{Theorem}}{}
\makeatother
\title[Countable separation property for associative algebras]{Countable separation property for associative algebras}
\author{Alexey Petukhov}
\date{Source: arXiv:2510.25455v1 (CC BY 4.0)}
\begin{document}
\maketitle

\noindent\emph{Source and licence.} Countable separation property for associative algebras. Version arXiv:2510.25455v1 (CC BY 4.0), distributed under the Creative Commons Attribution 4.0 International licence, as recorded in the arXiv licence metadata for this exact version and shown on the abstract page https://arxiv.org/abs/2510.25455v1. Full rights notice: licenses/arxiv-2510.25455-CC-BY-4.0.txt.\par\smallskip

\noindent\textit{Editorial scope.} Counted unit: the Lemma whose source label is \texttt{L:stab} in the article (source0.tex lines 151--152, source label \texttt{L:stab}), delivered together with the article's own complete original proof of it (source0.tex lines 153--188, one \texttt{proof} environment of 36 lines). No mathematics is written, repaired or re-derived: statement and proof are the article's own text line for line. Required context retained uncounted: T:stab (lines 52--55). Every definition, display and label of those lines is retained unchanged. Omitted from this package and not redistributed: 1--51, 56--150, 189--323. Nothing inside a retained excerpt is omitted or altered: every retained body is delivered exactly as the article's lines are written, so no omission receipt of any kind accompanies this package. Citations: the retained excerpts carry no citation command at all, so no bibliography entry of the article is transcribed and no reference is redistributed. Reference adaptation: none was needed and none was made; every reference inside the delivered unit resolves inside it, so no unresolved reference or citation is produced. Printed slips kept literal: no printed word, symbol, hypothesis or number of the counted statement or of its proof was altered. Layout and class: the article's own class is \texttt{amsart} with packages including xy, amssymb, amsmath, euscript, amsthm; the wrapper is the lane's pinned \texttt{amsart} editorial wrapper at 10pt on A4, which restates the theorem-like environments the retained text needs and adds the article's own macro definitions that the retained text needs (none), with extra wrapper packages (bm, bbm, dsfont, xspace, cancel, mathtools). The delivered numbering is the unit's own. Assets: no retained span contains a figure environment, a graphics-inclusion command, a PDF-page inclusion, a table or a TikZ picture; no image file is copied into this package, the package declares no asset dependency and no image file is redistributed with it. Licence: version arXiv:2510.25455v1 of this article is distributed under the Creative Commons Attribution 4.0 International licence, as recorded in the arXiv licence metadata for this exact version and shown on the abstract page https://arxiv.org/abs/2510.25455v1. A full rights notice is delivered at licenses/arxiv-2510.25455-CC-BY-4.0.txt. Characters: every retained excerpt is pure ASCII, so the delivered engine, XeLaTeX with its default Latin Modern text font, reports no missing character.
\begin{theorem}\label{T:stab} Let $A$ be an at most countable-dimensional associative algebra with 1 over a field $\mathbb K$ and let $\widetilde{\mathbb K}$ be an arbitrary extension of $\mathbb K$. \\
(a) if $A$ admits a simple module $M$ with the trivial annihilator such that ${\rm Hom}_{A}(M, M)=\mathbb K$ then $A$ is E-CSP.\\
(b) if $A\otimes_{\mathbb K}{\widetilde{\mathbb K}}$ admits a simple module $M$ with the trivial annihilator such that  ${\rm Hom}_{A\otimes_{\mathbb K}{\widetilde{\mathbb K}}}(M, M)=\mathbb K$ then $A$ is E-CSP.
\end{theorem}

\begin{lemma}\label{L:stab} Under the assumptions of Theorem~\ref{T:stab} let $(0)\ne V\subset A$ be a finite-dimensional $\mathbb K$-vector space of dimension $d$. 
Then there exists $a_1, \ldots, a_d\in A\backslash(0)$ such that if $I\in{\rm Var}(V)$ then $a_i\in I$ for some $i$.\end{lemma}

\begin{proof} 
Assume that there exists a counterexample to this lemma given by a finite-dimensional vector space $V$. 
Without loss of generality we can assume that $\dim_{\mathbb K} V$ is the least possible. 

If $d=\dim_{\mathbb K} V=1$ then every $v\in V\backslash(0)$ is contained in $I$ for all $I\in{\rm Var}(V)$.

We assume next that $d\ge2$. 
Fix $v\in V\backslash(0)$. 
For each $x\in A$ set
$$\phi_x: V\to A, \hspace{10pt} w\to wxv-vxw.$$
Assume that for some $x\in A, w\in V$ we have $\phi_x(w)\ne0$. 
Fix such $x$ and $w$. 
Put $V_1:={\rm Ker}\phi_x, V_2:={\rm Im}\phi_x$. 
Then $v\in V_1, 0\ne\phi_x(w)\in V_2$ and thus $V_1, V_2\ne (0)$. 
By definition we have $\dim_{\mathbb K} V_1+\dim_{\mathbb K} V_2=\dim_{\mathbb K} V$ and hence $\dim_{\mathbb K} V_1, \dim_{\mathbb K} V_2<\dim_{\mathbb K} V$. 
It is easy to see that if $I\in{\rm Var}(V)$ then either $I\in{\rm Var}(V_1)$ or $I\in{\rm Var}(V_2)$. 
By the assumption of the minimality of $\dim_{\mathbb K}V$ the statement of Lemma~\ref{L:stab} holds for $V_1, V_2$ providing the desired list $a_1,\ldots, a_d\in A\backslash(0)$.

We left to consider the case when $vxw=wxv$ for all $x\in A, w\in V$. 
Recall that $M$ is an exact simple $A$-module with ${\rm Hom}_{A}(M, M)=\mathbb K$. 
Pick $m\in M$ such that $vm\ne 0$ and pick $w\in V$ such that  $\mathbb Kw\ne\mathbb Kv$. 
Without loss of generality we assume that $wm\ne0$ --- if it is not the case we can replace $w$ by $w+v$. 

Assume $wm\notin\mathbb Kvm$. 
Density theorem implies that there exists $x\in A$ such that $xwm=0, xvm=m$ (here we use essentially that ${\rm Hom}_A(M, M)=\mathbb K$). 
Next we have $0=xwm=vxwm=wxvm=wm$. 
This contradicts our assumption $wm\ne0$. 

We left to consider the case $wm\in\mathbb Kvm$ or equivalently $wm=Cvm$ for some $C\in\mathbb K$. 
We claim that $wm'=Cvm'$ for all $m'\in M$. 
Indeed we have that for each $m'$ there exists $x\in A$ with $m'=xvm$ (here the simplicity of $M$ is essential). 
We have
$$wm'=wxvm=vxwm=Cvxvm=Cvm'$$
as claimed. 
Thus $(w-Cv)M=0$ and hence $M$ can't be an exact $A$-module.
\end{proof}

\end{document}
